\documentclass[11pt]{article}
\usepackage{finprobs}

\title{Does Unification Help? The Intermittency Parameter\\ and Joint SPX/VIX Calibration}
\author{finprobs Research Project}
\date{July 2026}

\begin{document}
\maketitle

\begin{abstract}
Pure rough Heston/Bergomi models are known to struggle with joint calibration to SPX and VIX option smiles simultaneously. We ask whether the log S-fBM model's extra intermittency parameter, $\lambda^2$, offers a structural fix. Working directly from the model's closed-form covariance, we prove that $\lambda^2$ cancels out of the ratio governing short-dated skew versus VIX-level vol-of-vol at fixed $H$, and is therefore calibration-inert for this purpose; the correlation-timescale parameter $T$ is the actual lever. We verify this cancellation numerically to machine precision. We then show, via a cross-paper derivation, that the same aggregation mechanism explaining index-versus-single-name roughness differences likely undermines rather than reinforces this $T$-lever at the index level -- a genuine tension between two of this project's most promising theoretical threads, not a synergy.
\end{abstract}

\section{Background and motivation}

\citet{gatheral2020quadratic} precisely characterize the joint SPX/VIX calibration failure of pure rough volatility models: matching the steep, near-divergent short-maturity SPX skew requires a level of vol-of-vol inconsistent with observed VIX-implied-volatility levels. This sits within the broader ``Guyon conjecture'' territory \citep{guyon2022vix} that continuous-price-path models may structurally struggle to fit both surfaces at once. Known fixes -- quadratic rough Heston \citep{gatheral2020quadratic}, path-dependent volatility \citep{guyonlekeufack2023pdv,gazzaniguyon2025fourfactor}, Gaussian polynomial Volterra models \citep{abijaber2025gaussian} -- all add structure beyond plain fBM-driven log-volatility.

\begin{finding}
No paper tests log S-fBM, or any structurally similar unified model, against the joint SPX/VIX calibration problem. We confirmed this independently via two separate literature searches. The source paper itself \citep{wu2022rough} is purely a probability/statistics paper with no options-pricing content; the authors explicitly flag ``a faithful model for asset and option prices'' as future work.
\end{finding}

\citet{abijaberilland2025joint} provide a relevant cautionary precedent: a conventional one-factor Markovian model \emph{outperforms} rough and path-dependent alternatives with the same parameter count on this exact benchmark. Extra structural sophistication does not automatically buy better joint fit.

\section{A derived negative result: $\lambda^2$ cancels out}

Recall the log S-fBM covariance (\citealp{wu2022rough}, verified directly from the source):
\begin{equation}
C(\tau;H,T,\nu^2) = \frac{\nu^2}{2}\left[T^{2H}-\tau^{2H}\right], \qquad \nu^2=\frac{\lambda^2}{H(1-2H)}. \label{eq:cov}
\end{equation}

\begin{theorem}
For fixed $H$ and $T$, the ratio $R(\tau_1,\tau_2) \equiv C(\tau_1;H,T,\nu^2)/C(\tau_2;H,T,\nu^2)$ is independent of $\lambda^2$ for all $\tau_1,\tau_2<T$.
\end{theorem}
\begin{proof}
$\nu^2$ enters \eqref{eq:cov} as a common multiplicative prefactor: $C(\tau;H,T,\nu^2)=\nu^2\cdot g(\tau;H,T)$ where $g(\tau;H,T)=\tfrac12[T^{2H}-\tau^{2H}]$ does not depend on $\nu^2$ (equivalently, on $\lambda^2$). Hence $R(\tau_1,\tau_2)=g(\tau_1;H,T)/g(\tau_2;H,T)$, in which $\nu^2$ -- and therefore $\lambda^2$ -- does not appear.
\end{proof}

Taking $\tau_1=\tau_{\text{short}}$ (a short SPX maturity) and $\tau_2=\tau_{\text{VIX}}$ ($\approx$ one month), Theorem~1 says the ratio between short-dated and VIX-horizon vol-of-vol is entirely controlled by $H$ and $T$, and \emph{not at all} by $\lambda^2$ -- for any fixed $H$, no choice of intermittency can differentially target short-dated skew steepness relative to the VIX level.

\begin{corollary}
$T$ is the parameter that governs $R(\tau_{\text{short}},\tau_{\text{VIX}})$: as $T\to\infty$ (the pure-fBM limit), $R\to1$ and the lever vanishes; as $T\to\tau_{\text{VIX}}^+$, $R\to\infty$, its maximal differentiating power.
\end{corollary}

This makes $T$ -- not $\lambda^2$ -- structurally analogous to the second, slower timescale already used in existing multi-factor fixes for the joint calibration problem \citep{gazzaniguyon2025fourfactor}.

\section{Numerical verification}

We evaluated \eqref{eq:cov} directly at $H=0.10$, $T=0.5$ years, $\tau_{\text{short}}=1/252$, $\tau_{\text{VIX}}=21/252$.

\begin{table}[h]
\centering
\begin{tabular}{lccccc}
\toprule
$\lambda^2$ & 0.01 & 0.05 & 0.10 & 0.50 & 2.00 \\
\midrule
$R(\tau_{\text{short}},\tau_{\text{VIX}})$ & 2.058205 & 2.058205 & 2.058205 & 2.058205 & 2.058205 \\
\bottomrule
\end{tabular}
\caption{The ratio is invariant to $\lambda^2$ to six decimal places across a $200\times$ range, confirming Theorem~1 exactly.}
\end{table}

\begin{table}[h]
\centering
\begin{tabular}{lrrrrrrrr}
\toprule
$T$ & 0.10 & 0.25 & 0.50 & 1.00 & 2.00 & 5.00 & 20.0 & $10^5$ \\
\midrule
$R$ & 13.28 & 2.86 & 2.06 & 1.71 & 1.51 & 1.36 & 1.23 & 1.03 \\
\bottomrule
\end{tabular}
\caption{$R(\tau_{\text{short}},\tau_{\text{VIX}})$ is strictly monotone decreasing in $T$, confirming Corollary~1.}
\end{table}

\begin{finding}
Theorem~1 and Corollary~1 are now verified facts, not merely derived claims: the cancellation is exact to numerical precision, and the $T$-dependence is exactly as predicted in both limiting regimes.
\end{finding}

\section{A cross-paper tension: does the $T$-lever survive aggregation?}

The companion paper on aggregation (this series, Paper~2) shows that cross-sectional averaging over heterogeneous single-name log-vol processes pulls the index's Hurst exponent toward that of a smoother common factor once $\beta^4N_s\gg10$. We ask whether the same mechanism affects the index's effective $T$.

\begin{proposition}
Under the additive factor model $\omega_i(t)=\beta_i F(t)+X_i(t)$ with independent components, the index covariance is
\[
\Cov(\Omega) = \Big(\textstyle\sum_i w_i\beta_i\Big)^2 C_F(\tau) + \Big(\textstyle\sum_i w_i^2\Big) C_X(\tau).
\]
As $N_s$ grows, $\sum_i w_i^2\to0$ (equal-weighted case, at rate $N_s^{-1}$ or faster under the diversification scaling of Paper~2), so $\Cov(\Omega)\to(\sum_iw_i\beta_i)^2 C_F(\tau)$ pointwise in $\tau$ -- the \emph{entire} covariance profile, truncation range included, converges to the common factor's shape, not merely its short-lag behavior.
\end{proposition}

If $T_F > T_{X}$ -- plausible if market-wide volatility regimes persist longer than firm-specific idiosyncratic shocks, though this is an economic premise we do not establish -- then the index's effective $T$ is pulled toward the common factor's (longer) value under the same mechanism that raises its effective $H$.

\begin{finding}
\label{fnd:tension}
Combining this with Corollary~1: $R(\tau_{\text{short}},\tau_{\text{VIX}})$ is strictly decreasing in $T$, so a longer effective index-level $T$ \emph{weakens} the very lever needed to decouple SPX skew from VIX level. Since SPX is itself an aggregate over roughly 500 names, the diversification mechanism that would make it a well-behaved index (Paper~2) simultaneously erodes the mechanism proposed here to fix its joint calibration problem. This is a genuine, non-forced tension between two of this project's most promising theoretical threads -- not a synergy -- and it generates a new falsifiable prediction: $T_{\text{index}}$ should open up specifically once $\beta^4N_s\gtrsim10$, the same threshold governing $H_{\text{index}}$.
\end{finding}

\section{Findings summary}

\begin{itemize}
\item \labelestablished. The joint SPX/VIX calibration problem is real and precisely characterized \citep{gatheral2020quadratic}.
\item \labelnegative. $\lambda^2$ cannot be the fix (Theorem~1), verified numerically to machine precision.
\item \labelplausible. $T$ is a structurally sensible lever (Corollary~1), analogous to existing multi-factor fixes, but untested against real option data.
\item \labelplausible. The aggregation-driven erosion of the $T$-lever at index level (Finding~\ref{fnd:tension}), conditional on the unestablished premise $T_F>T_X$.
\item \labelspeculative{}/gap. Whether unification helps joint calibration at all remains untested; no pricing engine for log S-fBM currently exists.
\end{itemize}

\section{Future directions}

\begin{enumerate}
\item Build a minimal log S-fBM option-pricing/calibration engine (likely simulation-based, since the model is not obviously affine or Markovian) and test the $T$-lever hypothesis directly against public SPX/VIX option-chain data.
\item Include an explicit parameter-matched overfitting control, following the precedent of \citet{abijaberilland2025joint}, so that any apparent improvement cannot be attributed merely to an extra free parameter.
\item Test the new prediction from Finding~\ref{fnd:tension} directly: estimate $T$ separately for single names and the index across sub-portfolios of varying $N_s$ and mean $\beta$, checking whether $T_{\text{index}}$ opens up at the same threshold that governs $H_{\text{index}}$.
\item If the premise $T_F>T_X$ fails empirically, reassess whether a genuinely index-level (not inherited via aggregation) fast timescale could be sourced separately -- e.g.\ from market-wide liquidity or leverage-driven vol-of-vol -- to preserve the $T$-lever's usefulness despite the aggregation tension.
\end{enumerate}

\section*{Acknowledgments}
This paper synthesizes research conducted by an orchestrated multi-agent AI investigation, with all citations independently re-verified against primary sources and all numerical claims computed and verified before inclusion here.

\bibliographystyle{plainnat}
\begin{thebibliography}{99}

\bibitem[Abi~Jaber and Li(2025)]{abijaber2025gaussian}
E.~Abi~Jaber and S.~Li.
\newblock Volatility models in practice: rough, path-dependent, or Markovian?
\newblock \emph{Mathematical Finance}, 2025. arXiv:2401.03345.

\bibitem[Abi~Jaber et al.(2025)]{abijaberilland2025joint}
E.~Abi~Jaber, C.~Illand, and S.~Li.
\newblock Joint SPX-VIX calibration with Gaussian polynomial volatility models.
\newblock \emph{Mathematical Finance}, 2025. arXiv:2212.08297.

\bibitem[Gatheral et al.(2020)]{gatheral2020quadratic}
J.~Gatheral, P.~Jusselin, and M.~Rosenbaum.
\newblock The quadratic rough Heston model and the joint S\&P~500/VIX smile calibration problem.
\newblock \emph{arXiv:2001.01789}, 2020.

\bibitem[Gazzani and Guyon(2025)]{gazzaniguyon2025fourfactor}
G.~Gazzani and J.~Guyon.
\newblock Pricing and calibration in the 4-factor path-dependent volatility model.
\newblock \emph{Quantitative Finance}, 25(3), 2025. arXiv:2406.02319.

\bibitem[Guyon(2022)]{guyon2022vix}
J.~Guyon.
\newblock The VIX future in Bergomi models: essential properties and calibration.
\newblock \emph{SIAM Journal on Financial Mathematics}, 13(4), 2022.

\bibitem[Guyon and Lekeufack(2023)]{guyonlekeufack2023pdv}
J.~Guyon and J.~Lekeufack.
\newblock Volatility is (mostly) path-dependent.
\newblock \emph{SSRN 4174589}, 2023.

\bibitem[Wu et al.(2022)]{wu2022rough}
Q.~Wu, J.-F.~Muzy, and E.~Bacry.
\newblock From rough to multifractal volatility: the log S-fBM model.
\newblock \emph{Physica A}, 604:127919, 2022. arXiv:2201.09516.

\end{thebibliography}

\end{document}
